% TOP_PROBLEM: {"id": 30006916, "title": "Martin's Conjecture for Definable Turing-Invariant Functions", "category_id": 18, "category": {"id": 18, "name": "logic", "display_name": "Logic", "description": "Problems in mathematical logic, model theory, proof theory, and finite model theory.", "slug": "logic", "order_index": 18, "created_at": "2026-07-31T15:26:25.671Z"}, "status": "open"}
% FIRST_POSTED: 2026-09-25
% ENTRY_KIND: statement_only
% !TeX program = lualatex
% Definition and sources only; this file is not an LLM solution attempt.
\documentclass[11pt]{article}
\usepackage[margin=25mm]{geometry}
\usepackage{fontspec}
\setmainfont{Latin Modern Roman}
\usepackage{amsmath,amssymb,mathtools}
\usepackage{unicode-math}
\setmathfont{Latin Modern Math}
\usepackage{xurl,hyperref}
\hypersetup{colorlinks=true,urlcolor=blue}
\setcounter{secnumdepth}{0}
\setlength{\parindent}{0pt}
\setlength{\parskip}{5pt}
\setlength{\emergencystretch}{3em}
\begin{document}
\section*{\texorpdfstring{Martin's Conjecture for Definable Turing-Invariant Functions}{Martin's Conjecture for Definable Turing-Invariant Functions}}\label{30006916}
\subsection{Definitions and mathematical statement}
Assume the Axiom of Determinacy in the source's ZF setting. For subsets of ${\mathbb{N}}$, $X\le_TY$ means $X$ is computable with oracle $Y$, and $X'$ is the oracle halting problem. A cone is $\{X:X\ge_TA\}$ for some $A$. A function $f:\mathcal P({\mathbb{N}})\to\mathcal P({\mathbb{N}})$ is Turing-invariant if $X\equiv_TY$ implies $f(X)\equiv_Tf(Y)$. Put $f\le_Mg$ when $f(X)\le_Tg(X)$ throughout some cone. Martin's two assertions are: every such $f$ is constant in Turing degree on a cone or satisfies $X\le_Tf(X)$ on a cone; and the functions in the latter class, modulo cone equivalence, are well-ordered by $\le_M$, with successor $[f]\mapsto[X\mapsto f(X)']$. ``Increasing'' here means above the identity, not necessarily order-preserving. The jump-iterate slogan is expressed by this precise successor law, without inventing a convention for every transfinite iterate.
\par\smallskip\textit{Formulation and concept references:} \href{https://arxiv.org/html/2510.19147}{[S1]}.

\subsection{Short English statement}
Under determinacy, are Turing-invariant functions eventually constant or above the identity, with the nonconstant cone-equivalence classes well-ordered and successive classes obtained by the Turing jump?
\subsection{Sources}
\begin{itemize}
\item[S1]Antonio Nakid Cordero, Martin's Conjecture in the Enumeration Degrees (2025), Introduction, Conjecture 1.2 and following paragraph.
\url{https://arxiv.org/html/2510.19147}.
\textit{Formulation source. Consulted 24 September 2026: Martin conjecture Parts I and II, cone order, and jump successors.}
\end{itemize}
\end{document}
